\chapter{Elliptische Kurven Kryptographie}
\label{sec:kryptographie}
Wir betrachten eine möglichkeit einen geheimen Schlüssel zu etablieren nach Diffie und Hellman. (zitiert nach Taylor und Francis 2008, S.170f.)
\begin{enumerate}
    \item Alice und Bob einigen sich auf eine elliptische Kurve $E$ über einem endlichen Körper $\mathbb{F}_q$, sodass das diskrete Logarithmusproblem in $E(\mathbb{F}_q)$ schwer zu lösen ist.
    \item Alice wählt ein geheimes $a$, das ein Integer sein soll, berechne $P_a=aP$ und schickt $P_a$ zu Bob.
    \item BoB wählt einen geheimen Integer b, er berechnet $P_b=bP$ und schickt Alice $P_b$.
    \item Alice berechnet $aP_b=abP$
    \item Bob berechnet $bP_a=baP$
    \item Alice und Bob benutzen eine öffentliche Methode um einen Schlüssel von $abP$ zu extrahieren. Zum Beispiel könnten sie die letzten 256 bits der x-Koordinate von $abP$ als Schlüssel verwenden oder sie könnten eine Hash-Funktion auf der x-Koordinate definieren.
\end{enumerate}
Die einzige Information die ein Außenstehender sieht, wäre die Kurve $E$, den Körper $F_q$ und die Punkte $P, aP,bP$. Daher muss sie das \textbf{Diffie-Hellman Problem} lösen.

\section{Warum eliptische Kurven verwenden?}
Der beste Algorithm um ECDLP zu lösen ist exponential, darum werden Elliptische Kurven Gruppen verwendet. (An Introduction to the Theory of Elliptic Curves, S.6)
Genauer der beste Algorithmus um ECDLP zu lösen über $\mathbb{F}_p$ braucht $O(\sqrt p)$ Zeit.
\section{Diffie-Hellman}
\label{sec:diffie-hellman}


\section{Sicherheitbetrachtung}
Faktorisierung großer Zahlen (Computeralgebra Kapitel 5, ab S.167)